% Problem-section file following ../_template.tex.
\section{Honest-party lockability of distillable key}
\label{sec:problem-22}

\paragraph{Problem.}
Can loss of one qubit held by an honest party reduce the two-way distillable
secret key by an arbitrarily large amount?  Let $K_D(A:B)_\rho$ denote the
asymptotic secret-key rate obtainable from $\rho_{AB}$ by local operations and
public two-way classical communication, with an adversary holding a
purification.  The question is whether there are finite-dimensional states
$\rho^{(r)}_{A_ra:B_r}$, with $\dim a=2$, such that
\begin{equation}
  K_D(A_ra:B_r)_{\rho^{(r)}}
  -K_D(A_r:B_r)_{\operatorname{Tr}_a\rho^{(r)}}
  \xrightarrow[r\to\infty]{}\infty.
  \label{eq:p22-ab-key-locking}
\end{equation}
Equation~\eqref{eq:p22-ab-key-locking} is the $AB$-locking question: the
discarded qubit belongs to Alice rather than being transferred to the
adversary.

\subsection*{Status}

Unsolved

\subsection*{Source}

Christandl et al. explicitly ask whether distillable key is lockable when the
discarded subsystem is held by an honest party
\sourcecite{ref:p22-unifying-key}{CEH+07}.

\subsection*{Progress}

\begin{itemize}
  \item Christandl and collaborators proved that transferring a bounded
  subsystem to the adversary cannot cause an unbounded loss of distillable
  key.  They explicitly left the distinct honest-party, or $AB$-, locking
  problem open, even for classical tripartite states
  \sourcecite{ref:p22-unifying-key}{CEH+07}.

  \item Horodecki and collaborators established non-lockability bounds for
  restricted families of private states.  They identify the arbitrary-state
  version of Eq.~\eqref{eq:p22-ab-key-locking} as unresolved, including a
  specified product-system special case
  \sourcecite{ref:p22-private-leakage}{HSK+21}.
\end{itemize}

\subsection*{References}

\begin{enumerate}[label={},leftmargin=4.8em,labelsep=0.5em]
  \footnotesize
  \item[\textup{[CEH+07]}]\label{ref:p22-unifying-key}
  M. Christandl, A. Ekert, M. Horodecki, P. Horodecki, J. Oppenheim, and
  R. Renner, ``Unifying Classical and Quantum Key Distillation,'' in
  \emph{Theory of Cryptography Conference 2007}, Lecture Notes in Computer
  Science \textbf{4392}, 456--478 (Springer, 2007).
  \href{https://doi.org/10.1007/978-3-540-70936-7_25}{doi:10.1007/978-3-540-70936-7\_25};
  \href{https://arxiv.org/abs/quant-ph/0608199}{arXiv:quant-ph/0608199}.

  \item[\textup{[HSK+21]}]\label{ref:p22-private-leakage}
  K. Horodecki, M. Studziński, R. P. Kostecki, O. Sakarya, and D. Yang,
  ``Upper Bounds on the Leakage of Private Data and an Operational Approach
  to Markovianity,'' \emph{Physical Review A} \textbf{104}, 052422 (2021).
  \href{https://doi.org/10.1103/PhysRevA.104.052422}{doi:10.1103/PhysRevA.104.052422};
  \href{https://arxiv.org/abs/2107.10737}{arXiv:2107.10737}.
\end{enumerate}

\subsection*{Comment}

The unresolved scope is $AB$-lockability of $K_D$ for arbitrary states, as
formalized in Eq.~\eqref{eq:p22-ab-key-locking}.  It excludes both the
non-lockable adversary-side operation and the restricted private-state
families covered by existing bounds.

\subsection*{Tag}

Quantum cryptography; Entanglement measures;
Local operations and classical communication; Resource conversion;
Qubit systems

\subsection*{ID}

\texttt{op\_f1e5a1cad168b38b}
